%%%%%%%%%%%%%%%%%%%%%%%%%%%%%%%%%%%%%%%%%%%%%%%%%%%%%%%%%%%%%%%%%%%%%%%%%%%%%%%
% CHAPTER 5: COMPLEMENTARITY AS A STRUCTURAL PRINCIPLE
%%%%%%%%%%%%%%%%%%%%%%%%%%%%%%%%%%%%%%%%%%%%%%%%%%%%%%%%%%%%%%%%%%%%%%%%%%%%%%%
% Version: 2.0 (January 2026)
%
% Purpose: Introduces complementarity as the central organizing principle of
%          EBF. Establishes supermodularity, coherence, and interaction
%          structures. Answers the third CORE question: "How do components interact?"
%
% Primary Appendix: B - CORE-HOW (Complementarity Structure)
% Secondary Appendices: X (Milgrom-Roberts), AC (Industrial Organization),
%                       AB (Matching Theory), BBB (Parameter Estimation)
%
% Prerequisites: Ch. 3 (Utility Theory), Ch. 4 (Agent Structure)
% Leads to: Ch. 6 (Reference Structures), Ch. 9 (Context)
%
% Chapter Type: A (CORE Chapter - answers HOW question)
%%%%%%%%%%%%%%%%%%%%%%%%%%%%%%%%%%%%%%%%%%%%%%%%%%%%%%%%%%%%%%%%%%%%%%%%%%%%%%%

\chapter{Complementarity as a Structural Principle}
\label{ch:complementarity}

%%%%%%%%%%%%%%%%%%%%%%%%%%%%%%%%%%%%%%%%%%%%%%%%%%%%%%%%%%%%%%%%%%%%%%%%%%%%%%%
% CHAPTER SCOPE BOX
%%%%%%%%%%%%%%%%%%%%%%%%%%%%%%%%%%%%%%%%%%%%%%%%%%%%%%%%%%%%%%%%%%%%%%%%%%%%%%%
\begin{tcolorbox}[
    colback=purple!5!white,
    colframe=purple!75!black,
    title={\textbf{Chapter Scope: Ziel / In-Scope / Out-of-Scope}},
    fonttitle=\bfseries
]

\textbf{Ziel (Objective):}
Establish the formal foundation for gamma interactions within the 10C CORE framework.

\vspace{0.3cm}
\textbf{In-Scope:}
\begin{itemize}[nosep]
    \item Core definitions and axioms for gamma interactions
    \item Integration with other 10C dimensions
    \item Formal mathematical foundations
    \item Key theorems and their implications
\end{itemize}

\vspace{0.3cm}
\textbf{Out-of-Scope (delegiert an Appendices):}
\begin{itemize}[nosep]
    \item Detailed proofs and derivations $\rightarrow$ CORE Appendix
    \item Empirical calibration $\rightarrow$ Appendix UNMAPPED_BBB
    \item Domain-specific applications $\rightarrow$ Application chapters
\end{itemize}

\end{tcolorbox}



%%%%%%%%%%%%%%%%%%%%%%%%%%%%%%%%%%%%%%%%%%%%%%%%%%%%%%%%%%%%%%%%%%%%%%%%%%%%%%%
% QUICK REFERENCE BOX
%%%%%%%%%%%%%%%%%%%%%%%%%%%%%%%%%%%%%%%%%%%%%%%%%%%%%%%%%%%%%%%%%%%%%%%%%%%%%%%
\begin{tcolorbox}[
    colback=gray!5!white,
    colframe=gray!75!black,
    title={\textbf{Begriffe in diesem Kapitel}},
    fonttitle=\bfseries
]
\small
\begin{tabular}{@{}ll@{}}
\textbf{Key Concept} & \textbf{Definition} \\
\midrule
gamma interactions & See Section \ref{sec:ch5-intro} \\
\end{tabular}
\end{tcolorbox}



%%%%%%%%%%%%%%%%%%%%%%%%%%%%%%%%%%%%%%%%%%%%%%%%%%%%%%%%%%%%%%%%%%%%%%%%%%%%%%%
% INTUITION BOX
%%%%%%%%%%%%%%%%%%%%%%%%%%%%%%%%%%%%%%%%%%%%%%%%%%%%%%%%%%%%%%%%%%%%%%%%%%%%%%%
\begin{tcolorbox}[
    colback=blue!5!white,
    colframe=blue!75!black,
    title={\textbf{Intuition: A Motivating Example}},
    fonttitle=\bfseries
]
Consider \textbf{Anna}, a professional facing a decision that involves gamma interactions.

She initially evaluates the choice using standard criteria, but something feels incomplete.
\textbf{Thomas}, her colleague, suggests considering additional dimensions beyond the obvious.

This chapter explores what Anna discovers when she applies the EBF framework---how
gamma interactions interact with broader welfare considerations.

\medskip
\textit{By the end of this chapter, you will understand why Anna's initial analysis
was incomplete and how the EBF approach provides a more comprehensive view.}
\end{tcolorbox}



%==============================================================================
% DIE 9 FRAGEN NAVIGATION BOX
%==============================================================================
\BCMQuestionNav{5}

%==============================================================================
% 10C CORE CONNECTION BOX
%==============================================================================
\begin{tcolorbox}[
    colback=blue!5!white,
    colframe=blue!75!black,
    title={\textbf{10C CORE Connection: HOW (Question 3)}},
    fonttitle=\bfseries
]
\textbf{The 3rd CORE Question:} How do utility components interact?

\medskip
\begin{tabular}{ll}
\textbf{CORE Appendix:} & B (CORE-HOW) \\
\textbf{Output:} & Complementarity matrix $\Gamma = [\gamma_{dd'}]$ \\
\textbf{Key Insight:} & Non-separability creates multiple equilibria \\
\textbf{Novel Contribution:} & Context-dependent complementarity $\gamma(\Psi)$
\end{tabular}

\medskip
\textbf{The Complementarity Matrices:}
\begin{center}
\small
$\Gamma$ (utility interactions) $\cdot$ $\Lambda$ (context interactions) $\cdot$ $\delta$ (modulation tensor)
\end{center}

\textit{Complementarity is not just mathematics---it determines whether systems have unique or multiple equilibria, whether change is gradual or discontinuous, and whether interventions succeed or fail.}
\end{tcolorbox}

%==============================================================================
% METHODOLOGICAL FRAMING: FEHR-COMPATIBILITY
%==============================================================================
\begin{tcolorbox}[
    colback=red!5!white,
    colframe=red!75!black,
    title={\textbf{Methodological Note: The Exclusion Principle}},
    fonttitle=\bfseries
]
\textbf{Following Appendix UNMAPPED_FRM (LIT-FEHR-METHOD):}

This chapter treats complementarity coefficients $\gamma_{ij}$ as \textbf{exclusion parameters}:
\begin{itemize}[nosep]
    \item \textbf{Null Hypothesis:} $H_0: \gamma_{ij} = 0$ (additivity, classical economics)
    \item \textbf{Alternative:} $H_1: \gamma_{ij} \neq 0$ (non-additivity, complementarity claim)
    \item \textbf{Scientific Value:} A complementarity claim gains value precisely where it can be \textit{refuted}
\end{itemize}

\medskip
\textbf{What this means:}
\begin{itemize}[nosep]
    \item Classical economics (Arrow-Debreu) is not ``wrong''---it is the \textit{null model}
    \item $\gamma \neq 0$ claims require empirical justification in specific contexts
    \item EBF does not ``integrate'' models---it identifies \textit{where} the null model fails
\end{itemize}

\textit{Complementarity is not a default assumption but a risky hypothesis that can fail.}
\end{tcolorbox}

%==============================================================================
% APPENDIX REFERENCES BOX
%==============================================================================
\begin{tcolorbox}[
    colback=yellow!10!white,
    colframe=orange!75!black,
    title={\textbf{Appendix References for This Chapter}},
    fonttitle=\bfseries
]
\begin{tabular}{lll}
\textbf{Appendix} & \textbf{Content} & \textbf{Section} \\
\midrule
B (CORE-HOW) & Complete $\Gamma$-matrix (15 parameters), UNMAPPED_CMP-11, UNMAPPED_CMP-12 & All sections \\
D (FORMAL-CSTAR) & Reference structure $C^*$ proofs & \ref{sec:comp-coherence} \\
X (LIT-MILGROM) & Supermodularity foundations & \ref{sec:comp-formal} \\
AC (DOMAIN-IO) & Firm-level complementarities & \ref{sec:comp-org} \\
AB (DOMAIN-MATCHING) & Market design applications & \ref{sec:comp-market} \\
BBB (CORE-WHERE) & $\gamma$ calibration values & \ref{sec:comp-calibration} \\
\textbf{RB (LIT-BENABOU)} & \textbf{Narrative modulation $\gamma(N)$, sign reversal} & \textbf{\ref{sec:comp-intro}} \\
\textbf{XVIII (LIT-FEHR-METHOD)} & \textbf{Exclusion Principle, $\gamma = 0$ null} & \textbf{Methodology} \\
G (REF-GLOSSARY) & Symbol definitions & All sections
\end{tabular}

\medskip
\textbf{Reading Guide:} This chapter provides conceptual foundations and intuition;
Appendix UNMAPPED_HOW contains the complete mathematical specification with all 163 parameters.
\end{tcolorbox}


%%%%%%%%%%%%%%%%%%%%%%%%%%%%%%%%%%%%%%%%%%%%%%%%%%%%%%%%%%%%%%%%%%%%%%%%%%%%%%%
%                          PART I: INTRODUCTION
%%%%%%%%%%%%%%%%%%%%%%%%%%%%%%%%%%%%%%%%%%%%%%%%%%%%%%%%%%%%%%%%%%%%%%%%%%%%%%%

\section{Introduction: Why Choices Interact}
\label{sec:comp-intro}

%------------------------------------------------------------------------------

The preceding chapters established \emph{who} has utility (Ch.~3) and \emph{what}
utility consists of (Ch.~4). This chapter addresses the crucial next question:
\emph{how} do these utility components interact?

\begin{intuition}[The Training-Technology Example]
Consider a firm choosing two policies: investment in employee training ($x_1$)
and adoption of advanced technology ($x_2$).

\textbf{Independent choices} ($\gamma_{12} = 0$):
The return to training doesn't depend on technology level.
Each choice can be optimized separately.

\textbf{Complementary choices} ($\gamma_{12} > 0$):
Trained employees get more out of advanced technology.
Advanced technology makes training more valuable.
The \emph{whole} exceeds the sum of \emph{parts}.

\textbf{Substitute choices} ($\gamma_{12} < 0$):
If we invest in training, we don't need technology (people improvise).
If we invest in technology, we don't need training (machines work).
One investment crowds out the other.
\end{intuition}

%------------------------------------------------------------------------------
% CENTRAL QUESTION BOX
%------------------------------------------------------------------------------
\paragraph{The Central Question.}
This chapter addresses the 3rd of the 9 BCM Questions:

\begin{center}
\fcolorbox{blue!75!black}{blue!5!white}{\parbox{0.85\textwidth}{
\textbf{Question 3 (HOW):} How do utility components interact?
Are they complements, substitutes, or independent?
}}
\end{center}

\begin{definition}[Complementarity]
Two choices $x_i$ and $x_j$ are \textbf{complements} when doing more of one
increases the return to doing more of the other:
\begin{equation}
\gamma_{ij} = \frac{\partial^2 U}{\partial x_i \partial x_j} > 0
\label{eq:complementarity-def}
\end{equation}
They are \textbf{substitutes} when $\gamma_{ij} < 0$ and \textbf{independent}
when $\gamma_{ij} = 0$.
\end{definition}

%------------------------------------------------------------------------------
% NARRATIVE MODULATION (UNMAPPED_CMP-11)
%------------------------------------------------------------------------------
\begin{tcolorbox}[
    colback=green!5!white,
    colframe=green!60!black,
    title={\textbf{Advanced: Narrative Modulation of $\gamma$ (Axiom UNMAPPED_CMP-11)}},
    fonttitle=\bfseries
]
\textbf{Key Insight from Bénabou, Falk \& Tirole (2018):}

The \textit{sign} of $\gamma_{ij}$ can flip based on prevailing moral narratives:
\begin{equation}
\gamma_{ij}(N) = \begin{cases}
\gamma^{+} > 0 & \text{if } N = N_R \text{ (responsibilizing narrative)} \\
\gamma^{-} < 0 & \text{if } N = N_E \text{ (exculpatory narrative)}
\end{cases}
\end{equation}

\textbf{Example:} Whether social image ($U_S$) and financial incentives ($U_F$) are complements or substitutes depends on how society frames the action:
\begin{itemize}[nosep]
    \item \textit{Responsibilizing} ($N_R$): ``You \textit{could} help but chose money'' $\Rightarrow$ $\gamma > 0$
    \item \textit{Exculpatory} ($N_E$): ``Everyone does it, you're not special'' $\Rightarrow$ $\gamma < 0$
\end{itemize}

\medskip
\textit{Full formalization: Appendix UNMAPPED_B (CORE-HOW), Axiom UNMAPPED_CMP-11. Narrative theory: Appendix UNMAPPED_RB (LIT-BENABOU).}
\end{tcolorbox}

%------------------------------------------------------------------------------
% ASKING-OFFERING COMPLEMENTARITY (UNMAPPED_CMP-12)
%------------------------------------------------------------------------------
\begin{tcolorbox}[
    colback=blue!5!white,
    colframe=blue!60!black,
    title={\textbf{Advanced: Asking-Offering Complementarity (Axiom UNMAPPED_CMP-12)}},
    fonttitle=\bfseries
]
\textbf{Key Insight from Bénabou, Jaroszewicz \& Loewenstein (2022):}

In help-seeking interactions, asking and offering can be complements or substitutes depending on rejection costs:
\begin{equation}
\gamma(\text{offer}, \text{ask}) = \begin{cases}
\gamma^{+} > 0 & \text{if } \lambda\alpha \text{ low (rejection less costly)} \\
\gamma^{-} < 0 & \text{if } \lambda\alpha \text{ high (``waiting trap'')}
\end{cases}
\end{equation}

\textbf{The Fear of Asking mechanism:}
\begin{itemize}[nosep]
    \item \textit{Rejection reveals} low generosity $g < \tilde{g}$, threatening self-image
    \item \textit{Discouragement effect}: Non-offers reduce asking propensity ($w^*_N > w^*_O$)
    \item \textit{Waiting trap equilibrium}: Both parties inefficiently wait for the other
\end{itemize}

\textbf{Applications:} Welfare non-take-up, workplace help-seeking, healthcare avoidance.

\medskip
\textit{Full formalization: Appendix UNMAPPED_B (CORE-HOW), Axiom UNMAPPED_CMP-12. Fear of Asking: Appendix UNMAPPED_RB (LIT-BENABOU), Axiom UNMAPPED_RB-10. Theory: MS-IB-018.}
\end{tcolorbox}

%------------------------------------------------------------------------------
% CHAPTER OVERVIEW
%------------------------------------------------------------------------------
\paragraph{Chapter Overview.}
This chapter proceeds as follows:
\begin{itemize}
    \item Section \ref{sec:comp-formal}: Formal foundations (supermodularity)
    \item Section \ref{sec:comp-heritage}: Intellectual heritage
    \item Section \ref{sec:comp-types}: Types of complementarity
    \item Section \ref{sec:comp-levels}: Complementarity at each level
    \item Section \ref{sec:comp-coherence}: The coherence index $K$
    \item Section \ref{sec:comp-classical}: Integration with classical economics
    \item Section \ref{sec:comp-calibration}: From theory to calibration
    \item Section \ref{sec:comp-computational}: Computational enablement (2026)
    \item Section \ref{sec:comp-integration}: Integration with 10C framework
    \item Section \ref{sec:comp-summary}: Summary and transition
\end{itemize}


%%%%%%%%%%%%%%%%%%%%%%%%%%%%%%%%%%%%%%%%%%%%%%%%%%%%%%%%%%%%%%%%%%%%%%%%%%%%%%%
%                    PART II: FORMAL FOUNDATIONS
%%%%%%%%%%%%%%%%%%%%%%%%%%%%%%%%%%%%%%%%%%%%%%%%%%%%%%%%%%%%%%%%%%%%%%%%%%%%%%%

\section{Formal Foundations: Supermodularity}
\label{sec:comp-formal}

%------------------------------------------------------------------------------

\begin{definition}[Supermodularity]
A function $f: \mathbb{R}^n \to \mathbb{R}$ is \emph{supermodular} if for all $i \neq j$:
\begin{equation}
  \frac{\partial^2 f}{\partial x_i \partial x_j} \geq 0
  \label{eq:supermodularity}
\end{equation}
Equivalently, for all $x, y \in \mathbb{R}^n$:
\begin{equation}
  f(x \vee y) + f(x \wedge y) \geq f(x) + f(y)
\end{equation}
where $x \vee y$ is the component-wise maximum and $x \wedge y$ the component-wise minimum.
\end{definition}

\paragraph{Interpretation.}
Supermodularity means: \emph{Doing more of one thing increases the return to doing more of another.}
The cross-partial derivative is non-negative.
Actions are complements, not substitutes.

\paragraph{The Lattice Formulation.}
The $\vee / \wedge$ formulation shows supermodularity is about \emph{order}:
moving ``up'' in all dimensions together is better than moving up in some and down in others.
This connects complementarity to coordination: everyone moving together beats mixed movement.

\begin{cross-reference}
\textbf{Appendix UNMAPPED_HOW} provides the complete mathematical treatment of supermodularity,
including proofs of equilibrium existence, lattice structure, and comparative statics.
\textbf{Appendix UNMAPPED_FND} details the Milgrom-Roberts contributions that established these foundations.
\end{cross-reference}


%%%%%%%%%%%%%%%%%%%%%%%%%%%%%%%%%%%%%%%%%%%%%%%%%%%%%%%%%%%%%%%%%%%%%%%%%%%%%%%
%                    PART III: INTELLECTUAL HERITAGE
%%%%%%%%%%%%%%%%%%%%%%%%%%%%%%%%%%%%%%%%%%%%%%%%%%%%%%%%%%%%%%%%%%%%%%%%%%%%%%%

\section{Intellectual Heritage}
\label{sec:comp-heritage}

%------------------------------------------------------------------------------

\subsection{Edgeworth (1881): Origins}

\citet{edgeworth1881} first noted that goods can be complements:
the marginal utility of coffee increases with sugar.
This is consumer-level complementarity: $\partial^2 U / \partial c_{coffee} \partial c_{sugar} > 0$.

\subsection{Milgrom and Roberts (1990): The Foundations}

\citet{milgromroberts1990compare} established the mathematical foundations
for analyzing complementarity in economic systems. Their key results:
\begin{enumerate}
  \item Pure strategy Nash equilibria exist in supermodular games (no convexity needed)
  \item Equilibria form a complete lattice (smallest and largest equilibria exist)
  \item Comparative statics are monotone (higher parameters $\to$ higher equilibrium)
  \item Best responses are isotone (if others do more, optimal response is more)
\end{enumerate}

\subsection{Modern Manufacturing}

In a companion paper, \citet{milgromroberts1990modern} applied complementarity theory
to explain the transformation of manufacturing. They identified a cluster of
complementary practices:
\begin{itemize}[nosep]
  \item Shorter production runs + Reduced inventories (JIT)
  \item Improved quality control (TQM) + Flexible manufacturing
  \item Increased worker training + Flatter hierarchies
\end{itemize}

The key insight: these practices are \emph{complements}, not independent choices.
Adopting one raises the return to adopting others:
\begin{equation}
  \gamma_{JIT,TQM} > 0, \quad \gamma_{JIT,flex} > 0, \quad \gamma_{TQM,training} > 0
\end{equation}

This explains why firms adopt coherent ``bundles'' rather than mixing approaches.

\subsection{Empirical Evidence: HRM Bundles}

\citet{ichniowski1997} tested complementarity in human resource management.
Studying 36 steel production lines, they found:
\begin{itemize}[nosep]
  \item Individual HRM practices had small effects
  \item ``Bundles'' of complementary practices raised productivity by 7\%
  \item Piecemeal adoption was ineffective
\end{itemize}

This confirms the theoretical prediction: with strong complementarities,
coherent bundles dominate mixed strategies.

\subsection{Topkis (1998): The Mathematics}

\citet{topkis1998} provided rigorous mathematical foundations, establishing:
\begin{itemize}[nosep]
  \item Supermodular games have pure-strategy Nash equilibria
  \item Equilibria form a complete lattice
  \item Best responses are monotone
  \item Comparative statics are well-behaved
\end{itemize}


%%%%%%%%%%%%%%%%%%%%%%%%%%%%%%%%%%%%%%%%%%%%%%%%%%%%%%%%%%%%%%%%%%%%%%%%%%%%%%%
%                    PART IV: TYPES OF COMPLEMENTARITY
%%%%%%%%%%%%%%%%%%%%%%%%%%%%%%%%%%%%%%%%%%%%%%%%%%%%%%%%%%%%%%%%%%%%%%%%%%%%%%%

\section{Types of Complementarity}
\label{sec:comp-types}

%------------------------------------------------------------------------------

Complementarity manifests in several distinct forms:

\subsection{Strategic Complementarity}

\begin{definition}[Strategic Complementarity]
My best response increases when you do more:
\begin{equation}
  \frac{\partial^2 \Pi_i}{\partial a_i \partial a_j} > 0
  \quad \Rightarrow \quad
  \frac{\partial a_i^*(a_j)}{\partial a_j} > 0
\end{equation}
\end{definition}

\paragraph{Examples.}
\begin{itemize}[nosep]
  \item \textbf{Bank runs:} If others withdraw, I should withdraw (Diamond-Dybvig)
  \item \textbf{Technology adoption:} If others adopt, adoption is more valuable
  \item \textbf{Cooperation:} If others cooperate, cooperation is less risky
\end{itemize}

\paragraph{Consequence: Multiple Equilibria.}
Strategic complementarity generates multiple equilibria---high (all cooperate)
and low (none cooperate)---both self-fulfilling.

\subsection{Temporal Complementarity}

\begin{definition}[Temporal Complementarity]
My action today increases the return to my action tomorrow:
\begin{equation}
  \frac{\partial^2 U}{\partial a_t \partial a_{t+1}} > 0
\end{equation}
\end{definition}

\paragraph{Examples.}
\begin{itemize}[nosep]
  \item \textbf{Habit formation:} Exercise today makes exercise tomorrow easier
  \item \textbf{Human capital:} Learning today increases return to learning tomorrow
  \item \textbf{Identity:} Honest behavior today reinforces honest self-image
\end{itemize}

\paragraph{Consequence: Path Dependence.}
Temporal complementarity creates momentum: early choices shape later choices.

\subsection{Cross-Domain Complementarity}

\begin{definition}[Cross-Domain Complementarity]
My action in domain $A$ increases the return to action in domain $B$:
\begin{equation}
  \frac{\partial^2 U}{\partial a^A \partial a^B} > 0
\end{equation}
\end{definition}

\paragraph{Examples.}
\begin{itemize}[nosep]
  \item \textbf{Health-wealth:} Health enables earning; wealth enables health investment
  \item \textbf{Trust-trade:} Social trust enables exchange; trade builds trust
  \item \textbf{Education-income:} Skills enable earnings; income enables education
\end{itemize}

\paragraph{Consequence: Spillovers.}
Cross-domain complementarity means interventions in one domain affect others.

\subsection{Cross-Agent Complementarity}

\begin{definition}[Cross-Agent Complementarity]
My utility increases when you do more:
\begin{equation}
  \frac{\partial^2 U_i}{\partial x_i \partial x_j} > 0
\end{equation}
\end{definition}

\paragraph{Examples.}
\begin{itemize}[nosep]
  \item \textbf{Network effects:} Your participation makes the platform valuable to me
  \item \textbf{Herd immunity:} Your vaccination protects me
  \item \textbf{Social norms:} Your compliance makes compliance easier for me
\end{itemize}

\paragraph{Consequence: Coordination Problems.}
We could all be better off moving together, but individual incentives may not support it.


%%%%%%%%%%%%%%%%%%%%%%%%%%%%%%%%%%%%%%%%%%%%%%%%%%%%%%%%%%%%%%%%%%%%%%%%%%%%%%%
%                    PART V: COMPLEMENTARITY AT EACH LEVEL
%%%%%%%%%%%%%%%%%%%%%%%%%%%%%%%%%%%%%%%%%%%%%%%%%%%%%%%%%%%%%%%%%%%%%%%%%%%%%%%

\section{Complementarity at Each Level}
\label{sec:comp-levels}

%------------------------------------------------------------------------------

Complementarity operates at every level of the framework:

\subsection{Individual Level}
\label{sec:comp-individual}

\paragraph{Internal Coherence.}
An individual's choices across life domains are complementary:
\begin{itemize}[nosep]
  \item Values and behavior: Living according to values reinforces identity
  \item Health and productivity: Exercise improves work performance
  \item Learning and earning: Skills enable income; income enables learning
\end{itemize}

A person with $K_{individual} \approx 1$ has aligned goals, values, and actions.
Incoherence ($K < 1$) manifests as: inner conflict, procrastination, regret.

\subsection{Household Level}
\label{sec:comp-household}

\paragraph{Family Complementarity.}
Household members' choices interact:
\begin{itemize}[nosep]
  \item Specialization: Comparative advantage in earning vs. caring
  \item Coordination: Shared schedules, resources, goals
  \item Conflict: Zero-sum decisions, incompatible preferences
\end{itemize}

A family with $K_{household} \approx 1$ has clear roles, aligned priorities.
Incoherence manifests as: conflict, misunderstanding, divorce.

\subsection{Organization Level}
\label{sec:comp-org}

\paragraph{Organizational Complementarity.}
This is \citet{milgromroberts1990modern}'s core domain:
\begin{itemize}[nosep]
  \item Strategy and structure must fit
  \item Processes and culture must align
  \item Incentives and goals must match
\end{itemize}

A firm with $K_{org} \approx 1$ has all elements reinforcing strategy.
Incoherence manifests as: bureaucratic dysfunction, strategic failure.

\subsection{Market Level}
\label{sec:comp-market}

\paragraph{Market Complementarity.}
Markets exhibit complementarity patterns:
\begin{itemize}[nosep]
  \item Industry clusters: Firms benefit from proximity to related firms
  \item Platform economics: More users $\to$ more value $\to$ more users
  \item Supply chains: Specialized suppliers enable specialized producers
\end{itemize}

\subsection{National Level}
\label{sec:comp-national}

\paragraph{Institutional Complementarity.}
\citet{hallsoskice2001} documented varieties of capitalism:
\begin{itemize}[nosep]
  \item \textbf{Liberal (US, UK):} Flexible labor, dispersed ownership, general education
  \item \textbf{Coordinated (DE, JP):} Rigid labor, concentrated ownership, vocational training
\end{itemize}

Each is coherent ($K \approx 1$) but different. Mixing elements reduces coherence.

\subsection{Global Level}
\label{sec:comp-global}

\paragraph{International Complementarity.}
\begin{itemize}[nosep]
  \item Trade and peace: Economic interdependence raises cost of war
  \item Institutions and cooperation: WTO, IMF enable coordination
  \item Norms and behavior: Shared expectations stabilize interactions
\end{itemize}


%%%%%%%%%%%%%%%%%%%%%%%%%%%%%%%%%%%%%%%%%%%%%%%%%%%%%%%%%%%%%%%%%%%%%%%%%%%%%%%
%                    PART VI: THE COHERENCE INDEX
%%%%%%%%%%%%%%%%%%%%%%%%%%%%%%%%%%%%%%%%%%%%%%%%%%%%%%%%%%%%%%%%%%%%%%%%%%%%%%%

\section{The Coherence Index: Measuring Alignment}
\label{sec:comp-coherence}

%------------------------------------------------------------------------------

\begin{definition}[Coherence Index]
The coherence index measures alignment between actual and reference structure:
\begin{equation}
  K = 1 - \frac{\|C - C^*\|}{\|C^*\|}
\end{equation}
where $C$ is the observed complementarity matrix and $C^*$ is the self-consistent reference.
\end{definition}

\paragraph{Interpretation.}
\begin{itemize}[nosep]
  \item $K = 1$: Perfect alignment. Stable equilibrium.
  \item $K \approx 1$: Near alignment. System is stable.
  \item $K \ll 1$: Misalignment. Instability.
  \item $K \approx 0$: Chaos. Crisis.
\end{itemize}

\paragraph{Coherence vs. Quality.}
Coherence ($K$) is distinct from quality ($Q$):
\begin{itemize}[nosep]
  \item High $K$, high $Q$: Stable and good (Nordic countries)
  \item High $K$, low $Q$: Stable but bad (authoritarian equilibria)
  \item Low $K$, high $Q$ potential: Transitioning (reform periods)
  \item Low $K$, low $Q$: Unstable and bad (failed states)
\end{itemize}


%%%%%%%%%%%%%%%%%%%%%%%%%%%%%%%%%%%%%%%%%%%%%%%%%%%%%%%%%%%%%%%%%%%%%%%%%%%%%%%
%                    PART VII: CLASSICAL ECONOMICS
%%%%%%%%%%%%%%%%%%%%%%%%%%%%%%%%%%%%%%%%%%%%%%%%%%%%%%%%%%%%%%%%%%%%%%%%%%%%%%%

\section{Complementarity and Classical Economics}
\label{sec:comp-classical}

%------------------------------------------------------------------------------

\paragraph{The Hidden Assumption.}
Arrow-Debreu general equilibrium assumes $\gamma_{ij} = 0$:
\begin{equation}
  U_i(x_i, x_j) = U_i(x_i) \quad \Rightarrow \quad \frac{\partial^2 U_i}{\partial x_i \partial x_j} = 0
\end{equation}

Under this assumption:
\begin{itemize}[nosep]
  \item Each agent optimizes independently
  \item Market prices coordinate without explicit coordination
  \item Equilibrium is unique and efficient
\end{itemize}

\paragraph{What Happens When $\gamma_{ij} \neq 0$.}
With interdependence:
\begin{itemize}[nosep]
  \item Independent optimization is insufficient
  \item Prices may not coordinate effectively
  \item Multiple equilibria can exist
  \item Coordination becomes essential
\end{itemize}

\paragraph{The Framework's Contribution.}
We don't reject classical economics---we show it as a \emph{special case}.
When $\gamma_{ij} \approx 0$ (atomistic markets), classical results hold.
When $\gamma_{ij} \neq 0$ (social preferences, networks, institutions),
the fuller framework is needed.


%%%%%%%%%%%%%%%%%%%%%%%%%%%%%%%%%%%%%%%%%%%%%%%%%%%%%%%%%%%%%%%%%%%%%%%%%%%%%%%
%                    PART VIII: CALIBRATION
%%%%%%%%%%%%%%%%%%%%%%%%%%%%%%%%%%%%%%%%%%%%%%%%%%%%%%%%%%%%%%%%%%%%%%%%%%%%%%%

\section{From Theory to Calibration}
\label{sec:comp-calibration}

%------------------------------------------------------------------------------

\paragraph{The Calibration Imperative.}
Complementarity theory tells us \emph{that} $\gamma_{ij}$ matters and \emph{how} systems
behave when $\gamma_{ij} \neq 0$. But it does not tell us the \emph{magnitude}:
Is $\gamma_{training,technology} = 0.3$ or $0.8$? How does it vary with context?

This gap cannot be filled by theoretical reasoning alone. The strength of
complementarities is an \emph{empirical question} requiring calibration.

\paragraph{Why LLMs Cannot Fill This Gap.}
Large language models ``know'' about complementarity from texts.
But they cannot predict \emph{quantitative} strength because:
\begin{itemize}[nosep]
    \item \textbf{Text describes, data quantifies.} Literature rarely reports precise effect sizes.
    \item \textbf{Context-dependence is underspecified.} Does $\gamma_{JIT,TQM}$ differ between Japan and Germany?
    \item \textbf{Interaction effects are exponential.} With 8 $\Psi$-dimensions, there are 28 two-way interactions.
\end{itemize}

\paragraph{The Solution.}
Quantitative prediction requires calibrating $\gamma_{ij}(\Psi)$ on experimental data.
The 50 years of behavioral economics experiments provide this calibration base.

\begin{cross-reference}
\textbf{Appendix UNMAPPED_EST} (CORE-WHERE) provides the complete calibration framework,
including $\gamma$-values, $\delta$-coefficients, and epistemic status tags.
\end{cross-reference}


%%%%%%%%%%%%%%%%%%%%%%%%%%%%%%%%%%%%%%%%%%%%%%%%%%%%%%%%%%%%%%%%%%%%%%%%%%%%%%%
%                    PART IX: COMPUTATIONAL ENABLEMENT
%%%%%%%%%%%%%%%%%%%%%%%%%%%%%%%%%%%%%%%%%%%%%%%%%%%%%%%%%%%%%%%%%%%%%%%%%%%%%%%

\section{Why Complementarity Becomes Operational in 2026}
\label{sec:comp-computational}

%------------------------------------------------------------------------------

\paragraph{The Historical Barrier.}
Milgrom-Roberts wrote about complementarity in 1990.
Topkis formalized supermodularity in 1998.
Why did it take 25+ years to build a comprehensive framework?

The barrier was \emph{computational}. Consider what analyzing complementarity requires:

\textbf{Matrix Operations.}
The complementarity matrix $\Gamma \in \mathbb{R}^{n \times n}$ for $n$ choice variables
has $n^2$ entries. For our 6-dimension FEPSDE structure with context:
\begin{itemize}[nosep]
  \item $\Gamma$-matrix: 15 unique complementarity parameters
  \item $\Lambda$-matrix: 28 context interaction parameters
  \item $\delta$-tensor: 120 modulation parameters
  \item Total: 163 parameters requiring estimation
\end{itemize}

\paragraph{What Modern Infrastructure Provides.}

\textbf{1. GPU-Accelerated Linear Algebra.}
Matrix operations that took hours now take seconds.

\textbf{2. Automatic Differentiation.}
Computing $\partial C^* / \partial \Psi$ no longer requires hand-derived expressions.

\textbf{3. Fixed-Point Solvers.}
Anderson acceleration, continuation methods, homotopy algorithms.

\textbf{4. LLM-Assisted Implementation.}
Researchers can describe models in natural language and generate code automatically.

\begin{cross-reference}
\textbf{Appendix UNMAPPED_HRC} provides a detailed account of how computational capabilities
evolved from chalk-and-blackboard analysis in the 1960s to today's LLM-assisted frameworks.
\end{cross-reference}


%%%%%%%%%%%%%%%%%%%%%%%%%%%%%%%%%%%%%%%%%%%%%%%%%%%%%%%%%%%%%%%%%%%%%%%%%%%%%%%
%                    PART X: 10C INTEGRATION
%%%%%%%%%%%%%%%%%%%%%%%%%%%%%%%%%%%%%%%%%%%%%%%%%%%%%%%%%%%%%%%%%%%%%%%%%%%%%%%

\section{Integration with 10C Framework}
\label{sec:comp-integration}

%------------------------------------------------------------------------------

\subsection{The 10C Integration Table}

\begin{table}[htbp]
\centering
\caption{Chapter 5 Integration with 10C CORE Structure}
\label{tab:ch5-9c-integration}
\small
\renewcommand{\arraystretch}{1.3}
\begin{tabular}{@{}llp{5cm}l@{}}
\toprule
\textbf{CORE} & \textbf{Question} & \textbf{Connection to This Chapter} & \textbf{Reference} \\
\midrule
\cellcolor{purple!20}\textbf{WHO} & Who has utility? & Level-specific complementarity $\gamma^{(L)}$ & AAA \\
\cellcolor{red!20}\textbf{WHAT} & What is utility? & FEPSDE dimension interactions & C \\
\cellcolor{blue!20}\textbf{HOW} & How interact? & \textbf{This chapter: $\Gamma$-matrix} & B \\
\cellcolor{green!20}\textbf{WHEN} & Context? & Context-dependent $\gamma(\Psi)$ & V \\
\cellcolor{gray!20}\textbf{WHERE} & Parameters? & Calibrated $\gamma$-values & BBB \\
\cellcolor{orange!20}\textbf{AWARE} & Awareness? & Awareness of complementarities & AU \\
\cellcolor{violet!20}\textbf{READY} & Action threshold? & Complementarity in WAX formula & AV \\
\cellcolor{teal!20}\textbf{STAGE} & Journey phase? & Stage-specific complementarities & AW \\
\bottomrule
\end{tabular}
\end{table}

\subsection{Key Connections}

\paragraph{HOW $\times$ WHAT.}
The complementarity matrix $\Gamma$ operates on FEPSDE utility dimensions:
\begin{equation}
U_{total} = \sum_d w_d \cdot u_d + \sum_{d < d'} \gamma_{dd'} \cdot u_d \cdot u_{d'}
\end{equation}

\paragraph{HOW $\times$ WHEN.}
Complementarity is context-dependent:
\begin{equation}
\gamma_{dd'}(\Psi) = \bar{\gamma}_{dd'} + \sum_k \delta_{dd',k} \cdot (\Psi_k - 0.5)
\end{equation}

\paragraph{HOW $\times$ WHO.}
Different aggregation levels have different complementarity structures:
\begin{equation}
\Gamma^{(L)} = f(L, \Psi)
\end{equation}


%%%%%%%%%%%%%%%%%%%%%%%%%%%%%%%%%%%%%%%%%%%%%%%%%%%%%%%%%%%%%%%%%%%%%%%%%%%%%%%
%                    PART XI: SUMMARY
%%%%%%%%%%%%%%%%%%%%%%%%%%%%%%%%%%%%%%%%%%%%%%%%%%%%%%%%%%%%%%%%%%%%%%%%%%%%%%%

\section{Summary and Transition}
\label{sec:comp-summary}

%------------------------------------------------------------------------------

\begin{tcolorbox}[
    colback=blue!5!white,
    colframe=blue!75!black,
    title={\textbf{Chapter 5 Summary: Complementarity (CORE-HOW)}}
]
\textbf{Core Contribution:}
Establishes complementarity as the structural foundation of EBF.

\textbf{Key Concepts:}
\begin{itemize}[nosep]
    \item \textbf{Supermodularity:} $\partial^2 f / \partial x_i \partial x_j \geq 0$
    \item \textbf{Complementarity matrix:} $\Gamma = [\gamma_{dd'}]$
    \item \textbf{Coherence index:} $K = 1 - \|C - C^*\| / \|C^*\|$
    \item \textbf{Context modulation:} $\gamma(\Psi) = \bar{\gamma} + \sum_k \delta_k \Psi_k$
\end{itemize}

\textbf{Key Results:}
\begin{enumerate}[nosep]
    \item Complementarity creates multiple equilibria
    \item Coherent configurations are stable attractors
    \item Same logic applies from individuals to global systems
    \item Classical economics is the special case where $\gamma_{ij} = 0$
\end{enumerate}

\textbf{Transition:}
Having established \emph{how} components interact, the next chapters address:
\begin{itemize}[nosep]
    \item Chapter 6: Reference structures $C^*$
    \item Chapter 9: Context modulation (WHEN)
    \item Chapter 10: Welfare dimensions (WHAT)
\end{itemize}
\end{tcolorbox}

%------------------------------------------------------------------------------
% FORMAL DETAILS
%------------------------------------------------------------------------------
\subsection{Formal Details}
\label{sec:comp-formal}

The complete formal treatment appears in Appendix UNMAPPED_B (CORE-HOW). Key results:
\begin{itemize}[nosep]
    \item \textbf{Existence:} Under continuity and monotonicity, coherent equilibria exist
    \item \textbf{Uniqueness:} With $\gamma > \bar{\gamma}_{crit}$, unique stable configuration
    \item \textbf{Stability:} $\|C_t - C^*\| \leq k^t \|C_0 - C^*\|$ with $k < 1$ (exponential convergence)
\end{itemize}

%------------------------------------------------------------------------------
% READING PATH BOX
%------------------------------------------------------------------------------
\begin{tcolorbox}[
    colback=gray!10!white,
    colframe=gray!75!black,
    title={\textbf{Reading Paths from Chapter 5}}
]
\textbf{For Theorists:}
$\rightarrow$ Appendix UNMAPPED_HOW (complete $\Gamma$-matrix) $\rightarrow$ Appendix UNMAPPED_FND (Milgrom-Roberts)

\textbf{For Scientists:}
$\rightarrow$ Appendix UNMAPPED_MS (134 theories with $C_{ij}$ restrictions) $\rightarrow$ Appendix UNMAPPED_DER (formal proofs)

\textbf{For Practitioners:}
$\rightarrow$ Chapter 6 (reference structures) $\rightarrow$ Chapter 9 (context)

\textbf{For Empiricists:}
$\rightarrow$ Appendix UNMAPPED_EST (calibration) $\rightarrow$ Appendix LLM1 (LLM-MC estimation)

\textbf{For Full 10C Understanding:}
$\rightarrow$ Chapter 10 (WHAT) $\rightarrow$ Chapter 11 (AWARE) $\rightarrow$ Chapter 12 (READY)
\end{tcolorbox}

%%%%%%%%%%%%%%%%%%%%%%%%%%%%%%%%%%%%%%%%%%%%%%%%%%%%%%%%%%%%%%%%%%%%%%%%%%%%%%%
